% SHARD-ID: AQM-67-QSA-FINITE-SET-FIRST-ASSOCIATION
% SHARD-TITLE: Structureless Finite-Set First Association
% SHARD-SUMMARY: Defines the first structureless finite-set association as a formal complex carrier with basis labels indexed by the set.
% SHARD-SUMMARY: Works the empty, one-point, and n-point cases by local derivation from the recorded convention.
% SHARD-SUMMARY: Records that any inner product or Gram matrix is extra data, not determined by the bare finite set.
% SHARD-KEYWORDS: quantum-system association, finite set, formal carrier, C^X, Gram matrix, Hilbert caveat

\section{Structureless Finite-Set First Association}
\label{sec:qsa-finite-set-first-association}

\subsection{Status and Local Anchors}

This is QSA Step 03 from
\path{docs/quantum_system_associations/PLAN.md}.  It is a local derivation
from \texttt{CONVENTIONS.md} conventions \texttt{(ap)} and \texttt{(aq)},
together with the planning contract in AQM-64, the vocabulary rule in AQM-65,
and the finite-set rows in AQM-66.  It uses no external source and introduces
no theorem about canonical Hilbert-space inner products.

The result is deliberately narrow.  The input is a bare finite set \(X\), with
no order, topology, measure, field-valued function space, site Hilbert spaces,
symplectic label space, or dynamics.  The output is only a formal complex
carrier with basis labels indexed by \(X\).

\subsection{First Carrier}

Under convention \texttt{(aq)}, the first structureless finite-set association
is
\[
  X \longmapsto \mathbb C\langle X\rangle
  :=
  \bigoplus_{x\in X} \mathbb C e_x .
\]
The Step 03 shorthand \(\mathbb C^X\) means this same formal carrier in this
finite-set context.  The symbols \(e_x\) are basis labels, not vectors with a
preassigned length, norm, phase convention, or orthogonality relation.

An element is displayed as
\[
  v=\sum_{x\in X} a_x e_x,\qquad a_x\in\mathbb C.
\]
This is the whole association at this stage:
\[
  \mathsf{FSet}_1(X)
  =
  \bigl(\mathbb C\langle X\rangle,\{e_x\}_{x\in X}\bigr).
\]
It does not return an adjoint, a \(*\)-observable algebra, a Hamiltonian, a
tensor-product site system, a direct-sum particle interpretation, or
Weyl--Heisenberg data.

\subsection{Checked Edge Cases}

The checks in this subsection are just the finite bookkeeping consequences of
the displayed definition and convention \texttt{(aq)}.

\paragraph{Empty set.}
For \(X=\emptyset\),
\[
  \mathbb C\langle\emptyset\rangle
  =
  \bigoplus_{x\in\emptyset}\mathbb C e_x
  =
  0.
\]
There are no basis labels and no Gram entries.  Treating this zero carrier as
a Hilbert-space edge case, or comparing it with an empty tensor product, is not
part of this shard.

\paragraph{One-point set.}
For \(X=\{\ast\}\),
\[
  \mathbb C\langle\{\ast\}\rangle=\mathbb C e_\ast,
  \qquad
  v=a e_\ast .
\]
The association records one formal basis label.  A one-by-one Gram entry such
as \(G_{\ast\ast}\) would be an added choice; it is not supplied by the bare
point.

\paragraph{\(n\)-point set.}
For a displayed finite set \(X=\{x_1,\ldots,x_n\}\),
\[
  \mathbb C\langle X\rangle
  =
  \mathbb C e_{x_1}\oplus\cdots\oplus \mathbb C e_{x_n}.
\]
The listing \(x_1,\ldots,x_n\) is only a display choice.  It lets us write
coordinates \((a_1,\ldots,a_n)\) for
\(\sum_{i=1}^n a_i e_{x_i}\), but the ordered coordinate display is not extra
structure returned by the bare-set association.

\subsection{Gram-Matrix Caveat}

An inner product is not determined in this workflow.  If a later construction
supplies a Gram array \(G=(G_{xy})_{x,y\in X}\), then one may write
\[
  \langle e_x,e_y\rangle_G=G_{xy}
\]
as notation for that added data.  This shard does not derive \(G\) from \(X\),
does not select the identity matrix, and does not classify which arrays are
acceptable inner products.  The data for a metric version would be
\((X,G)\), not the bare set \(X\).

Consequently, this is a structureless finite-set workflow, not a canonical
Hilbert-space inner-product theorem.  Words such as orthonormal, unitary,
adjoint, and Hilbert completion remain unavailable here unless a later shard
records the missing Gram or Hilbert convention.

\subsection{Comparison Boundary}

AQM-66 lists the finite-set carrier, tensor-site, direct-sum, and all-map
field-label workflows as separate comparison targets.  This shard makes only
the first carrier workflow available.  It proves no agreement, reduction, or
equivalence with the tensor-site systems planned for AQM-68, the direct-sum
particle systems planned for AQM-69, or the finite symplectic field-label
review planned for AQM-71.

In particular, special comparisons using one-dimensional internal spaces in a
future tensor or direct-sum workflow remain planned checks.  They are not
proved here, even in the one-point or \(n\)-point cases.
